\documentclass[spanish,letterpaper,oneside]{article}

\def\documenttitle {Cuatro Ideas Relevantes y Comentario Crítico de Cynthia Luna Scott (2015)}
\def\documentsubtitle {Fundamentos para la enseñanza y el aprendizaje I · Sesión 1}
\def\documentsubject {Actividad 1 · Cuatro Ideas de Cynthia Luna Scott (2015)}
\def\documentauthor {Mtro. Martín Jonathan de la Cruz Muñoz}
\def\coursename {Fundamentos para la enseñanza y el aprendizaje I}
\def\coursecode {MAEM-SD26}

\def\universityname {Instituto de Investigación, Innovación y Estudios de Posgrado para la Educación del Estado de Nuevo León}
\def\universityfaculty {Maestría en Aprendizaje y Enseñanza de las Matemáticas}
\def\universitydepartment {Fundamentos para la enseñanza y el aprendizaje I}
\def\universitydepartmentimage {img/iiiepe/marca-iiiepe.png}
\def\universitydepartmentimagecfg {height=1.75cm}
\def\universitylocation {Monterrey, Nuevo León}
\def\studentfullname {Mtro. Martín Jonathan de la Cruz Muñoz}
\def\studentid {00874}
\def\studentemail {martin.delacruz@campus.iiiepe.edu.mx}
\def\programperiod {2do periodo, septiembre--diciembre de 2026 (SD26)}
\def\programmodality {Modalidad presencial}
\def\programshort {MAEM}

\def\authortable {
  \begin{tabular}{ll}
    \textbf{Alumno:} & \studentfullname \\
    Matrícula: & \studentid \\
    Correo institucional: & \studentemail \\
    Unidad: & \coursename \\
    Programa: & \programshort \\
    Periodo: & \programperiod \\
    & \\
    \multicolumn{2}{l}{Fecha de realización: 19 de noviembre de 2026} \\
    \multicolumn{2}{l}{\universitylocation}
  \end{tabular}
}

\input{template}

\usepackage{helvet}
\renewcommand{\familydefault}{\sfdefault}
\usepackage{array}
\usepackage{longtable}
\usepackage{pdflscape}
\usepackage{tikz}
\usepackage{pgfplots}
\pgfplotsset{compat=1.18}
\usepackage[most]{tcolorbox}
\usetikzlibrary{arrows.meta,positioning,calc,fit,backgrounds,shapes.geometric}

\definecolor{iiiepeNavy}{HTML}{163B5C}
\definecolor{iiiepeBlue}{HTML}{1F6F9F}
\definecolor{iiiepeAqua}{HTML}{20B7A8}
\definecolor{iiiepeOrange}{HTML}{F28C28}
\definecolor{iiiepePaper}{HTML}{F6F8F7}
\definecolor{iiiepeInk}{HTML}{1E2A32}
\definecolor{iiiepePale}{HTML}{EEF5F8}
\definecolor{iiiepeMint}{HTML}{E7F6F4}
\definecolor{iiiepeWarm}{HTML}{FFF3E6}
\def\portraittitlecolor {iiiepeNavy}
\setcitestyle{authoryear,open={(},close={)}}
\makeatletter
\renewcommand{\@biblabel}[1]{}
\makeatother
\setlength{\bibhang}{1.27cm}
\setlength{\bibsep}{6pt}
\AtBeginDocument{\renewcommand{\refname}{Referencias}}

\newcommand{\insertcoverwatermark}{%
  \AddToShipoutPictureBG*{%
    \begin{tikzpicture}[remember picture,overlay]
      \node[opacity=0.04,inner sep=0pt] at (current page.center) {%
        \includegraphics[width=0.65\paperwidth]{img/iiiepe/marca-iiiepe.png}%
      };
    \end{tikzpicture}%
  }%
}

\setlength{\parskip}{5pt}\setlength{\parindent}{0pt}
\begin{document}\insertcoverwatermark\templatePortrait\templatePagecfg\templateFinalcfg\fancyfoot{}\fancyfoot[L]{\scriptsize IIIEPE · MAEM · SD26}\fancyfoot[C]{\scriptsize\thepage}\fancyfoot[R]{\scriptsize Subproducto 1B · Cuatro Ideas Scott}\setstretch{1.3}\raggedbottom
\section{Contexto de lectura y postura central}
A partir de la lectura integral del documento \textit{El futuro del aprendizaje 1: ¿Por qué deben cambiar el contenido y los métodos de aprendizaje en el siglo XXI?} de \citet{scott2015futuro}, selecciono y analizo críticamente cuatro ideas centrales junto con su comentario situado en la enseñanza de las matemáticas en bachillerato técnico e ingeniería.

Mi postura es que la transformación educativa descrita por \citet{scott2015futuro} no se resuelve con modernización superficial ni equipamiento tecnológico aislado, sino revisando simultáneamente qué aprenden los jóvenes, cómo se organizan las interacciones en el aula y bajo qué criterios humanistas se construye ciudadanía crítica \citep{unesco2015replantear}.

\section{Cuatro ideas relevantes y comentarios analíticos (Scott, 2015)}

\begin{tcolorbox}[colback=iiiepePale,colframe=iiiepeNavy,title=\textbf{Idea 1 · No basta «corregir el sistema»: transformar contenido, pedagogía y evaluación},breakable]
\textbf{Idea seleccionada de \citet{scott2015futuro}:} Las orientaciones escolares heredadas privilegian el recuerdo memorístico de contenidos estáticos y exámenes estandarizados incompatibles con la sociedad del conocimiento. Recuperando a Prensky, Saavedra y Opfer, \citet{scott2015futuro} advierte que «modernizar» superficialmente fracasa si no se transforman currículo, método y evaluación para cultivar pensamiento complejo, creatividad y resolución de problemas auténticos.\\[4pt]
\textbf{Comentario analítico:} En grupos de CONALEP, CECyTE e ingeniería resulta común encontrar estudiantes que resuelven despejes mecánicos cuando la instrucción es explícita, pero no deciden qué magnitud modela una tolerancia de taller o un costo operativo. Por tanto, el modelo transmisivo produce fluidez aparente sin comprensión relacional \citep{sawyer2008optimizing}; urge sustituir listas repetitivas por problemas que exijan modelar, validar unidades y argumentar en equipo.
\end{tcolorbox}

\begin{tcolorbox}[colback=iiiepeMint,colframe=iiiepeAqua!80!black,title=\textbf{Idea 2 · Nuevo aprendiz conectado, Pedagogía 2.0 y ciudadanía digital crítica},breakable]
\textbf{Idea seleccionada de \citet{scott2015futuro}:} El estudiantado habita múltiples pantallas y redes esperando participación activa y autonomía; cuando la escuela ignora esa realidad aumenta la apatía y el abandono. No obstante, \citet{scott2015futuro} aclara que la Pedagogía 2.0 y los dispositivos digitales no mejoran automáticamente el aprendizaje si carecen de intencionalidad pedagógica orientadora.\\[4pt]
\textbf{Comentario analítico:} Considero erróneo equiparar agilidad con aplicaciones móviles con alfabetización matemática. Si un estudiante grafica ecuaciones en software sin interpretar la pendiente como razón de cambio o el rango físico real, la herramienta acelera el error conceptual. En consecuencia, mi mediación docente debe enseñar a auditar toda salida digital mediante cálculo de contraste, análisis dimensional y justificación escrita \citep{sep2023pensamiento}.
\end{tcolorbox}

\begin{tcolorbox}[colback=iiiepeWarm,colframe=iiiepeOrange!90!black,title=\textbf{Idea 3 · Brecha escuela--trabajo, aprendizaje permanente y visión humanista},breakable]
\textbf{Idea seleccionada de \citet{scott2015futuro}:} La automatización desplaza labores rutinarias e incrementa la demanda de juicio experto y adaptabilidad ante problemas inéditos. Retomando los informes Faure, \citet{delors1996educacion} y Redecker, \citet{scott2015futuro} defiende estrechar el puente entre escuela y trabajo junto con una visión humanista integral \citep{unesco2015replantear}.\\[4pt]
\textbf{Comentario analítico:} En plantas industriales y laboratorios técnicos nadie entrega al egresado una ecuación ya despejada; los desafíos aparecen como fallas reales, variaciones térmicas o restricciones de presupuesto \citep{conalep2025mecatronica}. Sin embargo, reducir las matemáticas a entrenamiento laboral estrecho tampoco basta: la formación humanista habilita criterio ético, solidaridad y capacidad para seguir aprendiendo toda la vida.
\end{tcolorbox}

\begin{tcolorbox}[colback=iiiepePale,colframe=iiiepeBlue,title=\textbf{Idea 4 · Ecosistemas de aprendizaje interconectados y estudiante productor},breakable]
\textbf{Idea seleccionada de \citet{scott2015futuro}:} Frente al aislamiento del salón de clases, la educación contemporánea necesita articular cinco entornos (escuela, hogar, trabajo, comunidad y espacios digitales) dentro de ecosistemas abiertos donde el alumnado pase de consumidor pasivo a creador de conocimiento situado.\\[4pt]
\textbf{Comentario analítico:} En nuestros planteles tecnológicos, conectar el aula de matemáticas con mediciones del taller mecánico, costos domésticos de energía y discusión colaborativa en equipo vuelve significativo el conocimiento abstracto \citep{star2015algebra}. Además, esta apertura solo resulta equitativa cuando provee andamiajes accesibles a quienes parten con menor capital cultural previo.
\end{tcolorbox}

\clearpage
\section{Conclusiones}
En conclusión, sostengo que las cuatro ideas seleccionadas de \citet{scott2015futuro} convergen en un mandato claro para la enseñanza matemática: alinear tareas auténticas, uso crítico de tecnologías, rigor profesional y cuidado humanista del derecho a aprender. Por ello, considero que este marco orienta decisivamente mi transición en la maestría hacia secuencias comprensivas evaluables con evidencias reales de razonamiento.

\clearpage
\begin{thebibliography}{99}
\bibitem[CONALEP(2025)]{conalep2025mecatronica} Colegio Nacional de Educación Profesional Técnica. (2025). \textit{Perfil de egreso: Profesional Técnico-Bachiller en Mecatrónica. Modelo Académico 2025}. CONALEP.
\bibitem[Delors et~al.(1996)]{delors1996educacion} Delors, J., et al. (1996). \textit{La educación encierra un tesoro}. Santillana/UNESCO.
\bibitem[Sawyer(2008)]{sawyer2008optimizing} Sawyer, R. K. (2008). Optimizing learning: Implications of learning sciences research. En OECD, \textit{Innovating to learn, learning to innovate} (pp. 45--65). OECD Publishing.
\bibitem[Scott(2015)]{scott2015futuro} Scott, C. L. (2015). \textit{El futuro del aprendizaje 1: ¿Por qué deben cambiar el contenido y los métodos de aprendizaje en el siglo XXI?} UNESCO.
\bibitem[SEP(2023)]{sep2023pensamiento} Secretaría de Educación Pública. (2023). \textit{Programa de estudios del recurso sociocognitivo Pensamiento Matemático I}. SEP.
\bibitem[Star et~al.(2015)]{star2015algebra} Star, J. R., et al. (2015). \textit{Teaching strategies for improving algebra knowledge}. IES.
\bibitem[UNESCO(2015)]{unesco2015replantear} UNESCO. (2015). \textit{Replantear la educación: ¿Hacia un bien común mundial?} UNESCO.
\end{thebibliography}

\end{document}
